\documentclass[11pt]{article}

% ===== 基本宏包：尽量保持简单 =====
\usepackage[margin=1in]{geometry}
\usepackage{graphicx}
\usepackage{float}
\usepackage{chngcntr}
\usepackage{booktabs}
\usepackage{multirow}
\usepackage{longtable}
\usepackage{array}
\usepackage{calc}
\usepackage{amsmath, amssymb}
\setlength{\LTcapwidth}{\textwidth}
\usepackage{xurl}
\usepackage{hyperref}
\usepackage[numbers,sort&compress]{natbib}
\usepackage{setspace}
\usepackage{fontspec}
\usepackage{xeCJK}
\providecommand{\tightlist}{\setlength{\itemsep}{0pt}\setlength{\parskip}{0pt}}

% ===== 图表编号风格：全文连续编号，使用 Figure / Table =====
\renewcommand{\figurename}{Figure}
\renewcommand{\tablename}{Table}

% ===== 字体与换行策略 =====
\defaultfontfeatures{Ligatures=TeX,Scale=MatchLowercase}
\IfFontExistsTF{Noto Serif}{\setmainfont{Noto Serif}}{\IfFontExistsTF{DejaVu Serif}{\setmainfont{DejaVu Serif}}{\setmainfont{Latin Modern Roman}}}
\IfFontExistsTF{Noto Sans}{\setsansfont{Noto Sans}}{\IfFontExistsTF{DejaVu Sans}{\setsansfont{DejaVu Sans}}{\setsansfont{Latin Modern Sans}}}
\IfFontExistsTF{Noto Sans Mono}{\setmonofont{Noto Sans Mono}}{\IfFontExistsTF{DejaVu Sans Mono}{\setmonofont{DejaVu Sans Mono}}{\setmonofont{Latin Modern Mono}}}
\IfFontExistsTF{Noto Serif CJK SC}{\setCJKmainfont[AutoFakeBold=3,AutoFakeSlant=.2]{Noto Serif CJK SC}}{\IfFontExistsTF{Source Han Serif SC}{\setCJKmainfont[AutoFakeBold=3,AutoFakeSlant=.2]{Source Han Serif SC}}{\IfFontExistsTF{AR PL UMing CN}{\setCJKmainfont[AutoFakeBold=3,AutoFakeSlant=.2]{AR PL UMing CN}}{\setCJKmainfont[AutoFakeBold=3,AutoFakeSlant=.2]{WenQuanYi Zen Hei}}}}
\IfFontExistsTF{Noto Sans CJK SC}{\setCJKsansfont[AutoFakeBold=3,AutoFakeSlant=.2]{Noto Sans CJK SC}}{\IfFontExistsTF{Source Han Sans SC}{\setCJKsansfont[AutoFakeBold=3,AutoFakeSlant=.2]{Source Han Sans SC}}{\IfFontExistsTF{Droid Sans Fallback}{\setCJKsansfont[AutoFakeBold=3,AutoFakeSlant=.2]{Droid Sans Fallback}}{\setCJKsansfont[AutoFakeBold=3,AutoFakeSlant=.2]{AR PL UMing CN}}}}
\IfFontExistsTF{Noto Sans Mono CJK SC}{\setCJKmonofont[AutoFakeBold=3]{Noto Sans Mono CJK SC}}{\IfFontExistsTF{Droid Sans Fallback}{\setCJKmonofont[AutoFakeBold=3]{Droid Sans Fallback}}{\setCJKmonofont[AutoFakeBold=3]{AR PL UMing CN}}}
\setlength{\emergencystretch}{6em}
\sloppy

% 可选：让行距略宽一些，便于审稿阅读
\onehalfspacing

\title{\textbf{Decathlon VOC Analyzer：面向商品图文证据与用户评论对齐的证据驱动多模态 VOC 分析系统}}

\author{{\small
\begin{tabular}{@{}cccc@{}}
\multicolumn{4}{c}{\parbox[t]{0.6200\textwidth}{\centering \textbf{Severin Ye} \\ \textit{School of Computer Science and Engineering} \\ \textit{Kyungpook National University} \\ Daegu, Republic of Korea \\ {\ttfamily\small 6severin9\allowbreak@gmail.\allowbreak{}com}}} \\
\multicolumn{4}{c}{} \\
\parbox[t]{0.2200\textwidth}{\centering \textbf{Dokeun Lee} \\ \textit{Department of Physics} \\ \textit{Kyungpook National University} \\ Daegu, Republic of Korea \\ {\ttfamily\small nsa08008\allowbreak@naver.\allowbreak{}com}} & \parbox[t]{0.2200\textwidth}{\centering \textbf{Seowan Jung} \\ \textit{School of Computer Science and Engineering} \\ \textit{Kyungpook National University} \\ Daegu, Republic of Korea \\ {\ttfamily\small swan041014\allowbreak@gmail.\allowbreak{}com}} & \parbox[t]{0.2200\textwidth}{\centering \textbf{HyunJun Jung} \\ \textit{School of Computer Science and Engineering} \\ \textit{Kyungpook National University} \\ Daegu, Republic of Korea \\ {\ttfamily\small beanbeansummon\allowbreak@naver.\allowbreak{}com}} & \parbox[t]{0.2200\textwidth}{\centering \textbf{Hye Jeon} \\ \textit{School of Computer Science and Engineering} \\ \textit{Kyungpook National University} \\ Daegu, Republic of Korea \\ {\ttfamily\small prajna0426\allowbreak@naver.\allowbreak{}com}} \\
\end{tabular}
}}

\date{} % Nature 系列投稿一般不需要显示日期

\begin{document}
\maketitle

% =========================
% Abstract
% =========================
\begin{abstract}
电商商品的用户之声（Voice of Customer, VOC）分析需要同时理解用户评论中的主观体验和商品页面中的客观证据。仅依赖评论文本的情感或方面分析，难以判断用户观点是否被商品文案、规格或图像支持；直接把评论和商品图文整体交给大语言模型生成总结，又容易产生缺乏证据链的黑盒结论。针对这一问题，本文提出 Decathlon VOC Analyzer，一个面向商品图文证据与用户评论对齐的证据驱动多模态 VOC 分析系统。

该方法将原始商品信息标准化为包含商品文本、商品图片、图片局部区域和评论记录的证据包；在评论侧执行低信息过滤、星级抽样、方面级结构化抽取和去重；再将每个评论方面规划为具有明确检索意图和模态路线的问题；随后在商品文本和商品图像索引中执行两阶段召回与重排；最终生成包含优势、缺点、争议、证据缺口、改进建议和 claim 级证据归因的结构化报告。为支持复现和误差分析，系统保留关键中间表示、运行记录和评估接口，使报告不再是单段不可分解的自然语言摘要。

本文的主要贡献包括：第一，提出一种``评论方面建模---问题规划---多模态证据检索---证据约束生成''的商品 VOC 分析框架；第二，给出以结构化中间表示为核心的可复现实现，使评论、问题、证据、报告和评估指标均可追踪；第三，设计面向系统验证的实验协议，当前实现通过 166 项自动化测试，并支持 Recall@K、MRR、NDCG、claim support rate 和 citation precision 等指标。本文定位为系统方法论文，重点论证系统架构、实现机制和评估接口，而不声称已经完成覆盖多品类的大规模人工标注 benchmark。

关键词：多模态 RAG；用户之声分析；方面级评论建模；证据约束生成；商品评论分析；证据归因
\end{abstract}

% =========================
% Introduction
% =========================
\section{引言}
用户评论分析是电商智能运营中的基础任务。传统方法通常围绕情感分类、主题聚类或方面级情感分析展开，能够概括用户对商品属性的正负态度，却较少回答一个更关键的问题：这些评论观点是否能够被商品本体证据解释、支持或反驳。对于商品运营和产品改进而言，单纯知道``用户抱怨容量小''并不足够，分析者还需要知道商品页面是否明确描述了尺寸、图片是否展示了空间结构、该问题是商品缺陷还是用户预期偏差。

大语言模型为商品评论总结提供了新的可能，但一步式总结同样存在明显风险。如果系统直接将商品文案、图片描述和评论整体输入模型，模型可能生成流畅但难以审查的结论。使用者很难判断某条建议来自哪些评论、哪些商品文本或哪张图片，也难以定位错误发生在评论理解、检索、证据解释还是最终生成阶段。因此，商品 VOC 分析不应被简化为黑盒摘要任务，而应被建模为一条从评论需求到商品证据、再到结构化报告的可追踪链路。

本文研究的问题是：如何在电商商品场景中，将用户评论中的主观体验显式对齐到商品文案和商品图像中的可核验证据，并生成可审查、可复现、可评估的 VOC 洞察。为此，本文提出 Decathlon VOC Analyzer。系统首先以商品为中心构建标准化证据包，保留商品文本、商品图片、图片局部区域和评论记录；随后将评论抽取为方面级对象；再通过问题规划层把每个方面转化为文本、图像或图文联合检索任务；最后通过多模态召回、重排、证据聚合和 claim 级归因生成结构化报告。

本文的贡献可以概括为三点。

第一，本文提出一种面向商品 VOC 的证据驱动多模态 RAG 框架。不同于直接用评论原句检索商品证据，本文在评论方面和商品证据之间加入问题规划层，使主观表达被转化为可执行的证据查询。

第二，本文实现一套 artifact-first 的系统原型。系统中的商品证据、评论方面、检索问题、召回记录、报告洞察、claim attribution、feedback、replay 和 manifest 均以结构化 schema 表示并可落盘复用。

第三，本文给出可复现实验接口和评估设计。该设计同时覆盖检索指标与生成证据支撑指标，从而能够分别衡量``系统是否找到了相关证据''和``生成结论是否被证据支持''。

本文其余部分组织如下：第 2 节给出背景和问题定义；第 3 节讨论相关工作；第 4 节详细介绍系统方法；第 5 节说明实验设置与实现环境；第 6 节报告系统验证结果；第 7 节讨论方法意义；第 8 节总结全文；第 9 节说明局限性。

% =========================
% 背景与问题定义
% =========================
\section{背景与问题定义}
\subsection{商品 VOC 分析的证据需求}

面向商品的 VOC 分析与通用评论摘要不同。商品分析通常需要输出稳定的洞察对象，包括商品优势、主要缺点、争议点、适用场景、证据缺口和改进建议。每个洞察都应尽可能回指到两类证据：一类是用户评论中的原始表达，另一类是商品页面中的文本、图片或局部视觉结构。

如果系统只处理评论，它只能回答``用户怎么说''；如果系统只处理商品页面，它又无法知道``用户真正关心什么''。因此，商品 VOC 分析的关键在于把评论侧的主观需求和商品侧的客观证据连接起来。本文将这一连接建模为方面驱动的证据检索问题，而不是直接生成总结。

\subsection{输入与输出}

给定一个商品数据包 \(P\)，其包含商品文本 \(P_{text}\)、商品图像 \(P_{image}\) 和评论集合 \(C_{reviews}\)。系统目标是输出结构化报告 \(R\)，其中包括报告主张、支持证据、证据缺口和改进建议。该过程可表示为：

\[
R = F(P_{text}, P_{image}, C_{reviews})
\]

但本文中的 \(F\) 不是单步生成函数，而是一条分阶段链路：

\[
C_{reviews} \rightarrow A_{aspects} \rightarrow I_{intents} \rightarrow Q_{questions} \rightarrow E_{retrieval} \rightarrow G_{aggregates} \rightarrow R_{attributed}
\]

其中，\(A_{aspects}\) 表示评论方面对象，\(I_{intents}\) 表示问题意图，\(Q_{questions}\) 表示可执行检索问题，\(E_{retrieval}\) 表示商品图文证据，\(G_{aggregates}\) 表示方面聚合统计，\(R_{attributed}\) 表示带证据归因的最终报告。

\subsection{设计约束}

本文系统建立在五项约束之上。

第一，商品证据必须对象化。商品标题、描述、类目、图片、图片区域和评论不能只作为拼接文本输入模型，而应拥有稳定 ID 和来源字段。

第二，评论侧必须先结构化。系统需要把评论拆解为 aspect、sentiment、opinion、evidence span、usage scene 和 confidence，才能支持后续统计与检索。

第三，检索查询必须问题化。评论原句往往口语化且噪声高，直接检索容易产生宽泛候选；问题规划能够将其转化为可核验的证据需求。

第四，文本证据与图像证据应分 route 管理。商品文案和商品图片支持不同类型的判断，系统需要保留 route 信息用于召回、重排和归因。

第五，最终报告必须可追踪。报告中的主张应能映射到评论 ID、文本块 ID、图片 ID 或区域 ID；证据不足时应输出 evidence gap，而不是强行确认结论。

\subsection{与通用 RAG 问答的差异}

通用 RAG 系统通常以用户提出的问题为查询入口，目标是生成单个答案。本文任务的查询入口来自评论方面，目标不是单一回答，而是商品 VOC 画像。系统需要同时处理多个方面、多个问题、多个证据路线和多个报告字段，并保留完整过程 trace。因此，本文更接近一个面向商品分析的证据编排系统，而不是一个普通问答系统。

% =========================
% 相关工作
% =========================
\section{相关工作}
\subsection{检索增强生成与证据约束}

检索增强生成（Retrieval-Augmented Generation, RAG）将外部检索与参数化生成分离，使生成模型能够基于显式证据回答知识密集问题 \citep{lewis2020retrievalAugmentedGeneration}。后续研究进一步表明，检索不仅可以作为生成前的上下文补充，也可以用于生成后的校正与修订，例如 RARR 通过外部搜索和修订降低模型输出中的不实陈述 \citep{gao2023rarr}。这些工作为本文提供了基本方法论：商品 VOC 报告不应完全依赖模型内部知识，而应围绕可检索商品证据生成。

本文与通用 RAG 的不同在于，检索查询不是用户即时问题，而是由评论方面规划得到的证据问题；检索对象也不是百科段落，而是商品文本、商品图片和图片区域。最终目标不是单条答案，而是带 claim 级证据归因的结构化商品分析报告。

\subsection{多模态检索与视觉证据}

CLIP 证明了图像和文本可以通过自然语言监督映射到共享语义空间，从而支持跨模态检索 \citep{radford2021clip}。这一思想为商品图片作为一等证据对象提供了基础。多模态 RAG 的工程实践也强调，真实业务数据通常由文本、图像、表格和结构化字段共同构成，系统需要同时处理模态表示、跨模态检索和上下文融合 \citep{nvidia2024multimodalRagIntro,nvidia2024aiReadyKnowledgeSystems}。

视觉文档检索和原始模态召回研究进一步指出，图像不应过早被压缩为纯文本。ColPali 等方法将文档页面作为视觉对象进行检索，并通过多向量表示和 late interaction 提升细粒度匹配能力 \citep{faysse2024colpali}。本文没有直接实现 ColPali 式多向量检索，但吸收了保留原始视觉证据和向区域级证据扩展的思想，在商品图片层面同时建模整图和默认局部区域。

\subsection{评论分析与方面建模}

方面级情感分析通常将评论拆解为属性、观点和情感极性，适合发现用户集中关注的商品特征 \citep{pontiki2014semevalAspect}。对于商品 VOC 分析而言，方面建模仍是必要入口，因为评论经常包含多种体验、场景和情绪。如果系统不先进行方面级建模，后续检索和报告生成就只能依赖冗长评论原文。

本文与传统方面级情感分析的差异在于，方面不是最终结果，而是证据检索的起点。每个方面都会被转化为一个或多个带检索路线的问题，并进一步连接商品文本和商品图像证据。

\subsection{两阶段检索、向量库与重排}

实际检索系统常采用粗召回和精排分离的多阶段架构。粗召回负责覆盖候选，精排负责在小候选集上使用更高成本模型提高排序质量。Qdrant 的多阶段检索实践和 multi-vector search 课程均强调，在复杂检索任务中，向量库应作为候选管理基础设施，而不应替代后续 reranker 和业务判断 \citep{qdrant2024pdfRetrievalAtScale,qdrant2024multiVectorCourse}。

本文采用这一思想，将 embedding、索引后端、候选池构建和 reranker 解耦。文本候选和图像候选分别处理，最终在报告层统一归因。

\subsection{结构化输出与工作流编排}

LLM 在结构化输出任务中可能出现字段缺失、类型错误和格式漂移，因此需要 schema 约束、JSON mode 或语法约束等机制提高稳定性 \citep{liu2024structuredOutputBenchmark,geng2023grammarConstrainedDecoding}。ReAct 等工作则强调推理和行动之间的显式交替，说明复杂任务更适合拆为多步过程而非单轮提示 \citep{yao2023react}。

本文系统沿用这一方向：评论抽取、问题生成、报告生成均以结构化 schema 为边界；端到端流程则由 LangGraph 状态机组织为概览、标准化、索引和分析等节点。

\subsection{本文定位}

本文不提出新的通用 embedding 模型或 reranker 算法，而是提出一个面向商品 VOC 的系统化整合框架。其核心创新在于把评论方面、问题规划、多模态证据检索、结构化报告和 claim 级归因组织为一个可复现闭环，使商品分析从黑盒总结转向可审查的证据工作流。

% =========================
% 方法
% =========================
\section{方法}
\subsection{总体框架}

本文提出的 Decathlon VOC Analyzer 将商品 VOC 分析形式化为一个证据驱动的多阶段推理过程。系统首先把商品页面中的文本、图像和评论组织为统一证据空间；随后在评论侧抽取方面级用户反馈；再通过问题规划将主观反馈转化为可执行的证据查询；最后在商品文本与商品图像之间执行多模态检索、重排和证据约束生成。与一步式大模型摘要不同，该框架将``用户表达了什么''``商品证据能否支持''``最终报告如何归因''拆分为可观察的中间环节。

本文将该框架抽象为四个逻辑层。第一层是商品证据层，负责构建商品文本、图片和局部视觉区域的结构化表示。第二层是 VOC 需求层，负责将评论转化为方面、观点、情感和场景。第三层是证据对齐层，负责把方面改写为具有明确模态路线的问题，并在商品证据空间中检索和重排候选。第四层是报告归因层，负责生成结构化洞察，并将报告主张映射回评论证据与商品证据。

\begin{figure}[H]
\centering
\includegraphics[width=0.92\textwidth]{/home/severin/Codelib/Decathlon\_VOC\_Analyzer/论文\_md形式/图片/图1.png}
\caption{Decathlon VOC Analyzer 的总体框架。系统以商品证据包和评论方面为中间表示，通过问题规划连接主观 VOC 信号与多模态商品证据，并在报告阶段进行 claim 级归因。}
\end{figure}

\subsection{商品证据表示}

商品页面通常同时包含标题、描述、类目、规格、图片和用户评论。如果这些信息在系统入口处被简单拼接为长文本，后续模型虽然能够生成摘要，却很难追踪某一结论的来源。为此，本文首先构建商品级证据包，将商品文本、商品图片、图片局部区域和评论记录表示为具有稳定标识的证据对象。

文本证据保留其来源字段，例如标题、描述或类目。图像证据不仅保留整图，还构造若干规则化局部区域，以便后续模型能够在没有人工标注框的情况下初步访问局部视觉信息。评论证据则保留评分、语言提示和原文片段，为后续抽样与方面抽取提供输入。该表示的目的不是在数据层完成全部语义理解，而是为后续检索、归因和评估建立统一对象边界。

\subsection{评论方面建模}

评论侧的核心任务是将自然语言反馈转化为方面级 VOC 单元。每个方面单元包含被评价的属性、情感倾向、观点描述、原文证据片段、使用场景和置信度。与整句级摘要相比，方面级表示能够同时支持统计聚合和证据检索：系统既可以计算某一属性的正负反馈分布，也可以围绕该属性生成后续证据查询。

为了避免有限评论预算被高星短评占据，本文在评论抽取前引入星级感知抽样。默认策略更重视低星评论，从而优先捕获产品改进中更有价值的问题反馈。同时，为保证系统在不同运行条件下可复现，方面抽取同时支持结构化 LLM 路径和启发式路径。前者提供更强语义理解能力，后者保证在外部模型不可用时仍可完成端到端验证。

\subsection{从方面到问题的规划机制}

本文的关键设计是引入``方面到问题''的中间规划层。评论方面往往仍然是主观表达，例如``容量小''``佩戴舒适''``设计吸引儿童''。若直接用这些短语或原始评论检索商品证据，查询通常过宽，且难以判断应检索文本还是图像。因此，系统为每个方面构造若干具有明确检索意图的问题。

这些问题覆盖三类证据需求：文本支持、视觉确认和跨模态解释。文本支持问题关注商品文案或规格是否直接支持评论观点；视觉确认问题关注商品图片是否呈现相关结构或外观；跨模态解释问题则尝试判断某一反馈更接近真实产品问题、页面呈现不足还是用户预期偏差。通过这种方式，评论中的主观体验被转化为可以由商品证据回答的检索任务。

\begin{figure}[H]
\centering
\includegraphics[width=0.92\textwidth]{/home/severin/Codelib/Decathlon\_VOC\_Analyzer/论文\_md形式/图片/图2.png}
\caption{评论方面到证据查询的规划机制。方面对象并不直接进入检索，而是先被分解为具有不同证据意图的问题，再分别触发文本路线、图像路线或图文联合路线。}
\end{figure}

\subsection{多模态召回与重排}

证据对齐层采用两阶段检索。第一阶段使用向量表示进行粗召回，以较低成本从商品文本和商品图像中获得候选证据。第二阶段使用重排模型对候选集合进行精排，以提高最终证据的相关性。文本证据和图像证据在召回阶段被分路线管理，但在重排和报告阶段被统一到同一商品证据空间中。

文本路线主要用于验证商品名称、类目、描述、规格和显式承诺。图像路线主要用于验证结构、外观、颜色、局部细节和页面呈现方式。对于图像候选，系统既可以对整图建模，也可以对规则化区域进行裁剪和编码。虽然这些区域尚不等价于语义 grounding，但它们为局部视觉证据进入 VOC 分析提供了低成本接口。

为避免某一路线或某一种语言的候选过度占据结果，召回后还会构建语言与路线均衡的候选池。最终证据选择不仅考虑相关性分数，也考虑是否覆盖问题所声明的证据路线。因此，一个跨模态问题应尽量同时获得文本证据和图像证据，而不是被单一路线候选完全替代。

\subsection{证据约束报告与 claim 归因}

报告生成阶段不把召回结果作为普通上下文直接交给模型总结，而是先对方面信号和检索证据进行聚合。聚合对象记录每个方面的出现频次、情感分布、代表评论、证据片段和检索覆盖情况。随后系统生成结构化报告，报告字段包括优势、缺点、争议点、适用场景、证据缺口和改进建议。

为了降低生成式报告的不可审查性，本文进一步引入 claim 级归因。报告中的主要主张会被映射到评论证据、商品文本证据、商品图像证据或二者的组合，并标记为支持、部分支持、未支持或冲突。若某一评论观点在当前商品证据中无法得到充分验证，系统倾向于输出证据缺口，而不是把该观点包装成确定结论。

\begin{figure}[H]
\centering
\includegraphics[width=0.92\textwidth]{/home/severin/Codelib/Decathlon\_VOC\_Analyzer/论文\_md形式/图片/图3.png}
\caption{证据归因与评估闭环。报告主张被映射到评论、商品文本和商品图像证据；评估模块同时衡量检索质量和生成主张的证据支撑质量。}
\end{figure}

\subsection{可复现性与可解释性设计}

本文方法的可复现性来自两个层面。第一，所有关键中间结果均被表示为结构化对象，包括商品证据、评论方面、问题、检索记录、报告主张和评估指标。第二，系统在运行过程中保留可复用的中间产物和缓存，使研究者可以固定前序阶段，只替换问题生成、检索或重排策略进行消融。

这种设计也提升了错误分析能力。当最终报告出现不可靠建议时，分析者可以沿链路回溯到评论抽样、方面抽取、问题规划、候选召回、重排排序、报告生成或归因校准中的具体阶段。相比黑盒式摘要，本文框架更适合用于人工审查、模型替换、检索消融和后续多品类评测。

% =========================
% 实验设置
% =========================
\section{实验设置}
\subsection{实验目的}

本文的实验目标是验证所提出的系统框架是否能够稳定地完成证据驱动 VOC 分析，而不是报告一个已经冻结的大规模 benchmark。具体而言，实验关注四个问题：第一，系统是否能够从原始商品数据生成完整的结构化分析报告；第二，中间结果是否足以支持人工审查和错误定位；第三，真实模型链路与降级链路是否能够被明确区分；第四，当前评估接口是否能够承接后续人工标注数据。

这种实验定位与本文的系统方法性质一致。本文并不把当前样例集合视为公开标准数据集，也不声称当前结果可以外推到全部商品类别。相反，本文强调建立一套可复现的实验协议，使后续研究能够在同一框架下扩充数据、标注证据并比较不同检索策略。

\subsection{实现环境}

系统实现基于 Python 生态，采用 Web API 与批处理工作流两种入口。结构化对象由强类型 schema 约束，多阶段流程由状态图组织，模型调用通过兼容式语言模型网关完成。检索层支持本地索引和向量数据库两类后端，图像证据可通过视觉-语言表征模型编码，候选证据可通过文本或多模态重排模型重新排序。

为了区分正式实验与开发验证，系统提供运行策略控制。当完整模型链路可用时，系统使用真实 embedding、图像编码和重排模型；当外部能力不可用且策略允许时，系统退回启发式路径以保证流程可运行。该设计避免了将``模型能力不足''和``流程设计失败''混淆在一起。

\subsection{数据与验证单元}

实验数据来自商品页面抓取结果，包含商品文本、商品图片和用户评论。系统首先将这些原始信息转化为统一证据包，再围绕单个商品执行完整分析。当前阶段的验证单元是单商品运行实例，而非跨商品汇总指标。每次运行都会产生方面、问题、检索、报告和评估相关的结构化记录，从而形成可审查样本。

后续若要进行严格定量比较，需要在此基础上构建人工标注集。标注对象至少应包括问题级相关证据、报告 claim 的证据支撑状态、图像证据是否真正可见以及建议是否合理。本文当前实验设置已经为这些标注提供了对象边界和评估入口。

\subsection{评价指标}

本文将评价指标分为两类。第一类是检索指标。当存在问题级相关证据标注时，系统可计算 Recall@1、Recall@3、Recall@5、MRR、NDCG@3 和 NDCG@5。这些指标是信息检索评估中的常用排序与召回指标 \citep{manning2008introductionIr}，适合评估问题规划和多模态检索是否找到了正确商品证据。

第二类是生成证据支撑指标。对于结构化报告，系统统计主张是否被评论、商品文本或商品图像支持，包括 claim support rate、claim grounded rate、citation precision、citation contradiction rate 和不同模态路线的贡献比例。该类指标用于回答``报告是否被证据支持''，与传统只评价文本流畅度的摘要指标不同。

\subsection{工程验证协议}

由于本文关注系统可复现性，自动化测试也是实验协议的一部分。测试覆盖数据标准化、评论建模、问题生成、索引后端、embedding 与重排、检索服务、报告生成、HTML 导出、manifest 评估、运行策略和工作流入口。当前实现包含 166 项自动化测试且全部通过。这一结果说明系统的主要接口、结构化中间表示和工作流断言在当前实现中保持一致。

% =========================
% 结果与分析
% =========================
\section{结果与分析}
\subsection{工作流完整性}

实验首先验证系统是否能够从原始商品信息形成完整分析链路。结果表明，当前实现已经能够依次完成商品证据标准化、评论方面建模、问题规划、多模态召回、候选重排、方面聚合、报告生成和证据归因。各阶段输出具有一致的结构化表示，因此同一运行实例可以被拆解为多个可检查的分析单元。

这一结果说明，本文方法并非仅停留在概念设计层面。评论侧对象、商品侧证据和报告侧主张之间已经建立了明确的数据流。即使在关闭外部模型的验证模式下，系统仍保持相同的对象边界，使开发者能够区分流程正确性与模型能力差异。

\subsection{结构化可审查性}

第二项结果来自产物结构本身。一次完整分析不仅包含最终自然语言报告，还包含评论方面、问题意图、检索问题、召回证据、检索质量、运行配置、方面聚合、过程追踪和报告归因等对象。这些对象使分析者能够追问一条报告结论的来源，而不是只能接受模型给出的整体摘要。

这种结构化可审查性带来直接的错误定位能力。如果某一建议缺乏依据，可以检查它对应的评论是否被正确抽取、问题是否过宽、文本或图像路线是否召回了有效证据、重排是否把关键证据排在前列，以及报告后处理是否正确识别了证据缺口。相比单步摘要，这种错误定位粒度更适合研究迭代和人工审核。

\subsection{问题规划的作用}

结果分析表明，问题规划层是连接评论与商品证据的关键结构。原始评论经常包含情绪、背景和隐含假设，直接检索可能导致候选证据分散。方面级问题将这些表达转化为更明确的证据需求，例如``商品文案是否明确支持该体验''``商品图片是否展示相关结构''或``图文证据是否解释该负面反馈''。

由于每次检索都保留来源方面、来源问题和期望证据路线，系统能够在问题粒度上进行分析。后续人工标注也可以直接围绕问题判断哪些文本块、图片或区域是相关证据。这为比较不同检索策略提供了自然的评估单位。

\subsection{多模态证据的互补性}

文本证据和图像证据在商品 VOC 分析中承担不同功能。文本证据适合验证商品页面是否明确声明某项规格、用途或属性；图像证据适合观察结构、外观、颜色和局部细节。本文系统将二者作为不同路线检索，再在报告阶段统一归因，因此能够表达``评论观点由文本支持''``由图像支持''``由二者共同支持''或``当前证据不足''等不同状态。

这种 route-aware 设计也避免了图像证据被过早文本化。对于外观、结构和页面呈现相关问题，图像候选可以直接进入视觉模型重排，而不是完全依赖图片文件名或商品文案代理。虽然当前区域证据仍为规则化裁剪，但已经证明了将局部视觉证据纳入 VOC 链路的可行性。

\subsection{评估接口的可扩展性}

当前评估模块已经同时支持检索质量和报告归因质量。若后续补充人工相关性标签，系统可直接计算 Recall@K、MRR 和 NDCG；即使没有标签，系统仍能统计报告主张的支撑率、引用精度、冲突率和模态贡献。这使本文框架能够从系统验证自然过渡到正式 benchmark。

更重要的是，评估接口与中间对象共享同一标识体系。人工标注者可以标注问题、文本块、图片、区域或报告 claim，而不需要重新定义数据结构。该特征降低了未来扩展多品类评测集的成本。

\subsection{当前结果边界}

当前结果证明了系统流程完整、产物可审查、测试基线稳定，但仍不能被解释为最终性能结论。本文尚未报告冻结多品类数据集上的统计显著性比较，也尚未完成不同检索策略、不同重排模型和不同视觉区域策略的系统消融。因此，本节结果应被理解为系统方法验证，为后续大规模实验提供基础。

% =========================
% 讨论
% =========================
\section{讨论}
\subsection{为什么问题驱动检索适合商品 VOC}

商品 VOC 分析的难点不只是总结评论，而是解释评论中的体验是否能被商品证据支持。问题驱动检索在评论和商品之间增加了一个中间语义层：评论方面描述用户关注点，检索问题描述需要商品证据回答的判断。这使系统既不会脱离用户真实反馈，也不会让生成模型绕过商品证据直接给出结论。

这一设计尤其适合电商场景。商品页面天然包含多模态证据，评论又天然包含主观性和噪声。将二者直接拼接给模型会牺牲可控性，而将评论方面转写为证据问题，则可以把主观反馈拆解为可检索、可评估、可归因的操作单元。

\subsection{Artifact-first 的工程价值}

本文系统强调 artifact-first，即关键中间状态应被结构化保存，而不是仅存在于一次性模型调用上下文中。这种设计的价值在于，研究者可以复查证据标准化、方面抽取、问题规划、检索、重排和归因等任一阶段，而不是只能阅读最终报告。

Artifact-first 也降低了论文实验成本。后续若要比较不同问题生成策略，可以复用标准化产物；若要比较不同 reranker，可以复用评论方面和问题；若要做人审，可以直接基于 HTML 和 manifest 组织标注。

\subsection{与一步式大模型总结的对比}

一步式大模型总结实现简单，但通常难以回答三个问题：某个结论来自哪条评论，是否被商品文本支持，是否被商品图片支持。本文方法通过结构化 schema 和 claim attribution 将这三个问题显式化。系统复杂度确实更高，但换来了可解释性、可复现性和可评估性。

对于研究型系统而言，这种取舍是必要的。商品 VOC 报告不仅要``看起来合理''，还要能够被运营人员、产品人员或研究人员追问证据来源。若无法追踪证据，报告很难进入实际决策流程。

\subsection{对后续研究的启示}

当前实现说明，商品多模态分析的核心不一定是更大的生成模型，而是更好的证据组织和查询规划。后续研究可以沿四个方向扩展：第一，引入检测或分割模型，将固定 crop 区域升级为语义区域；第二，构建多品类人工标注集，系统评估 Recall@K、MRR 和 NDCG；第三，比较评论直召回、方面召回和问题驱动召回；第四，将 feedback/replay 与人工审查界面结合，形成持续修正机制。

\subsection{论文写作层面的意义}

本文当前版本已经可以作为系统方法论文的主体。它明确了问题定义、系统架构、模块实现、评估接口和局限性。后续若补充更大规模实验，可直接在现有章节中加入定量对照，而不需要重写方法主线。

% =========================
% 结论
% =========================
\section{结论}
本文提出 Decathlon VOC Analyzer，一个面向商品图文证据与用户评论对齐的证据驱动多模态 VOC 分析系统。系统将原始商品数据标准化为可追踪证据包，将用户评论抽取为方面级 VOC 对象，再通过问题规划、多模态召回、重排、报告生成和 claim 级归因生成结构化商品分析结果。

本文的核心结论是：商品 VOC 分析不应被简化为一步式评论摘要，而应被组织为``评论方面建模---证据问题规划---商品图文检索---证据约束生成''的多阶段流程。当前实现已经形成可运行的研究原型，并在结构化中间表示、证据归因、回放机制和评估接口上保持一致。测试结果显示当前实现通过 166 项自动化测试，说明系统主链和关键模块在工程层面具备稳定性。

未来工作将重点补充多品类人工标注 benchmark，系统比较不同检索和重排策略，引入更细粒度视觉 grounding，并将人工反馈与 replay 机制进一步闭环。

% =========================
% 局限性
% =========================
\section{局限性}
第一，本文当前结果仍属于系统验证，而非完整 benchmark。虽然系统已经支持 manifest 评估和自动化测试，但尚未构建覆盖多品类、多商品和多证据类型的冻结人工标注集。因此，本文不能声称系统在总体性能上优于所有基线。

第二，检索策略尚未完成系统消融。当前实现已经具备文本 route、图像 route、问题驱动检索、reranker 和缓存机制，但仍需要在统一数据集上比较评论直召回、方面直召回、问题驱动召回、不同 embedding 后端和不同 reranker 后端。

第三，视觉证据粒度仍有限。当前系统支持整图和固定比例区域 crop，但这些区域不是由语义检测、分割或 grounding 模型产生。对于鞋底结构、缝线、镜片光学细节等强局部属性，固定区域可能不足以提供精确解释。

第四，LLM 输出仍依赖提示词与 schema 约束。尽管系统提供结构化输出、后处理校准和 heuristic 兜底，但在复杂评论或跨语言场景下，方面抽取和报告生成仍可能出现误解、遗漏或过度概括。

第五，反馈闭环仍处于工程接口阶段。feedback/replay sidecar 已经存在，但尚未完成大规模人工审查流程，也未量化 replay 对后续报告质量的提升。

第六，参考文献和实验报告仍需面向正式投稿进一步扩展。本文当前更强调系统方法和实现事实，后续投稿版本应补充更系统的相关工作比较、正式图表、人工标注协议和统计显著性分析。

% =========================
% 致谢
% =========================
\section{致谢}
本稿基于当前仓库中的设计文档、源码实现、脚本、测试与输出产物整理而成。后续正式版本可在此补充项目指导、数据来源支持和实验环境协助等致谢信息。

% =========================
% 附录
% =========================
\section{附录}
\subsection{A. 关键目录与职责}

\begin{longtable}[]{@{}
  >{\raggedright\arraybackslash}p{(\columnwidth - 2\tabcolsep) * \real{0.5000}}
  >{\raggedright\arraybackslash}p{(\columnwidth - 2\tabcolsep) * \real{0.5000}}@{}}
\toprule\noalign{}
\begin{minipage}[b]{\linewidth}\raggedright
目录
\end{minipage} & \begin{minipage}[b]{\linewidth}\raggedright
职责
\end{minipage} \\
\midrule\noalign{}
\endhead
\bottomrule\noalign{}
\endlastfoot
\texttt{01\_data/} & 原始商品数据与中文审核数据 \\
\texttt{02\_outputs/} & 标准化结果、方面结果、索引、报告、feedback、replay、HTML 与 manifest \\
\texttt{03\_configs/} & 评论抽样 profile 与运行策略配置 \\
\texttt{04\_scripts/} & 工作流、评估、导出、验证和清理脚本 \\
\texttt{05\_src/} & 核心源码实现 \\
\texttt{06\_tests/} & API、服务、脚本和工作流测试 \\
\texttt{0\_docs/03\_论文子模块文档/} & 基于当前源码整理的论文模块文档 \\
\texttt{论文\_md形式/} & 论文分章稿、参考文献和导出脚本 \\
\end{longtable}

\subsection{B. API 入口}

系统当前包含以下主要 API：

\begin{itemize}
\tightlist
\item
  \texttt{/health}
\item
  \texttt{/api/v1/dataset/overview}
\item
  \texttt{/api/v1/dataset/normalize}
\item
  \texttt{/api/v1/index/overview}
\item
  \texttt{/api/v1/index/build}
\item
  \texttt{/api/v1/reviews/extract}
\item
  \texttt{/api/v1/analysis/product}
\end{itemize}

这些接口覆盖数据集概览、商品标准化、索引构建、评论结构化抽取和单商品综合分析。

\subsection{C. 复现命令}

端到端运行：

\begin{verbatim}
.venv/bin/python 04_scripts/run_workflow.py --category backpack --product-id backpack_010 --output-format json --export-html --write-manifest
\end{verbatim}

离线运行：

\begin{verbatim}
.venv/bin/python 04_scripts/run_workflow.py --category backpack --product-id backpack_010 --no-llm --output-format json
\end{verbatim}

验证多模态运行态：

\begin{verbatim}
.venv/bin/python 04_scripts/validate_multimodal_runtime.py --category backpack --product-id backpack_010
\end{verbatim}

评估 manifest：

\begin{verbatim}
.venv/bin/python 04_scripts/evaluate_manifests.py 02_outputs/7_manifests
\end{verbatim}

运行测试：

\begin{verbatim}
.venv/bin/python -m pytest
\end{verbatim}

合并论文：

\begin{verbatim}
.venv/bin/python 论文_md形式/脚本/build_complete_paper.py
\end{verbatim}

导出 PDF：

\begin{verbatim}
.venv/bin/python 论文_md形式/脚本/pipeline.py
\end{verbatim}

\subsection{D. 主要结构化产物}

\begin{longtable}[]{@{}lll@{}}
\toprule\noalign{}
产物 & 默认位置 & 用途 \\
\midrule\noalign{}
\endhead
\bottomrule\noalign{}
\endlastfoot
标准化商品证据包 & \texttt{02\_outputs/1\_normalized/} & 后续索引和分析输入 \\
评论方面结果 & \texttt{02\_outputs/2\_aspects/} & 问题规划和聚合输入 \\
索引与检索缓存 & \texttt{02\_outputs/3\_indexes/} & 召回和重排复用 \\
分析报告 & \texttt{02\_outputs/4\_reports/} & 最终结构化 VOC 结果 \\
feedback sidecar & \texttt{02\_outputs/5\_feedback/} & 人工反馈入口 \\
replay sidecar & \texttt{02\_outputs/5\_replay/} & 后续分析回放 \\
HTML 报告 & \texttt{02\_outputs/6\_html/} & 人工审查界面 \\
run manifest & \texttt{02\_outputs/7\_manifests/} & 评估与复现实验入口 \\
\end{longtable}

\subsection{E. 当前论文稿适用边界}

本稿适用于系统设计说明、方法章节草稿、项目答辩和后续实验规划。不适用于直接作为完整 benchmark 论文终稿，因为仍缺少多品类冻结标注集、正式消融实验和大规模人工评测。

% =========================
% References
% =========================
\bibliographystyle{unsrtnat}
\bibliography{ref}

\end{document}
